% id: 43-08
% book: 베이직쎈 대수 · 02 로그
% section: 유형 024 상용로그의 활용; 관계식이 주어질 때
% figure: none
% answer: 
% answer_source: 
% variant_level: 0 (원본 전사)
% 이 파일은 build.mjs 가 items.json 에서 만든다. 고칠 때는 items.json 을 고친다.
\smash{\raisebox{1.5mm}{\dmkichul{베이직쎈 대수 43쪽 08번}}}
디지털 사진을 압축할 때 원본 사진과 압축한 사진의 다른 정도를 나타내는 지표인 최대 신호 대 잡음비를 $P$, 원본 사진과 압축한 사진의 평균 제곱 오차를 $E$라 하면 다음과 같은 관계식이 성립한다고 한다.\par\smallskip\dispeq{$P=20\log 255-10\log E\ (E>0)$}\par\smallskip\noindent 두 원본 사진 $\pt{A}$, $\pt{B}$를 압축했을 때, 최대 신호 대 잡음비를 각각 $P_{\pt{A}}$, $P_{\pt{B}}$라 하고, 평균 제곱 오차를 각각 $E_{\pt{A}}\mkern3mu (E_{\pt{A}}>0)$, $E_{\pt{B}}\mkern3mu (E_{\pt{B}}>0)$라 하자. $E_{\pt{B}}=1000E_{\pt{A}}$일 때, $P_{\pt{A}}-P_{\pt{B}}$의 값을 구하시오.
